% !TEX program = xelatex
\documentclass[UTF8, zihao=5, a4paper, fontset=fandol]{ctexart}
\usepackage[answerkey]{gaokao-exam}

% 各页正文（不含页脚，去除空白字符后按 UTF-8 编码）SHA-256 摘要
\SetPageHash{1}{98dcff5c975261275256513e223721395e38b0c676cdb167154baf0b556456f7}
\SetPageHash{2}{f2a3606ce0c9a7c374559e91de799fb7e2e8653e3c1b34b7935a83bb66cd5569}
\SetPageHash{3}{2e3d9f76e7d1a8c833947bec96f9455dab2dff34afabbe469e626a30d85bb89d}
\SetPageHash{4}{7dcd0ecf2779a37ea3ca19c543cba0f63476d7181126aa2abbd87eeb86793b6f}

\begin{document}

\GaokaoTitleBlock
  {2026年普通高等学校招生全国统一考试}
  {外语（GPT 语）参考答案及评分参考}
  {（科目代码：99\quad 适用卷型：新课标 I 卷）}
  [评分参考 · 评卷专用]

\vspace{1.5mm}
\noindent{\fontsize{9.0pt}{12.5pt}\selectfont\fangsong
\textbf{【摘要核验说明】}第 1--45 题标准答案按题号顺序连续编码（\texttt{ASCII}，无空格或换行），所得 \texttt{SHA-256} 为：\\
\hspace*{2em}{\fontsize{8.4pt}{11pt}\selectfont\texttt{0ad9b5b9fb328591d9b4f58fb91fa72df051844c1f369fde9a1d8865d9b91e5a}}\\
非选择题各节摘要按考生定稿逐字重算核对；设同等给分答案的（见第 46--55 题），命中任一合规答案组合之摘要均予采信。}\par
\vspace{1.2mm}

% ============================================================================
% 第一部分  听力
% ============================================================================
\GaokaoPart{第一部分}{听力（每小题 3 分，满分 30 分）}

\begin{center}
\renewcommand{\arraystretch}{1.35}
\begin{tabular}{|c|c|c|c|c|c|c|c|c|c|c|}
\hline
\textbf{题号} & 1 & 2 & 3 & 4 & 5 & 6 & 7 & 8 & 9 & 10 \\
\hline
\textbf{答案} & \textbf{B} & \textbf{A} & \textbf{C} & \textbf{B} & \textbf{A} & \textbf{C} & \textbf{B} & \textbf{C} & \textbf{A} & \textbf{B} \\
\hline
\end{tabular}
\end{center}

\begin{enumerate}[label=\textbf{\arabic*.}, leftmargin=2.2em, itemsep=2pt, topsep=3pt]
  \item \textbf{B}\quad 脚本、示例配置和测试均已修改，目前仍在工作区，尚未提交。
  \item \textbf{A}\quad 母亲挑了小梨。“留四个”是女士随后建议的数量。
  \item \textbf{C}\quad 接口签名仍接受旧参数，实际超时仍为三十秒。
  \item \textbf{B}\quad 断言前排序保留了列表项数；少一项时两边长度仍不相等。
  \item \textbf{A}\quad 第二个文件保存本次复制的校验记录，测试只读取配置。
  \item \textbf{C}\quad 首次“通过”仅指三个单表样例，真实附件随后才打开测试。
  \item \textbf{B}\quad 两班表头逐字相同，可排除表头差异。
  \item \textbf{C}\quad 双进程抢锁用例依赖 Redis，首次运行时该服务未启动。
  \item \textbf{A}\quad 男士启动了 Redis 后重跑，测试用例未改动。
  \item \textbf{B}\quad 两个进程均记为取得同一把锁且中间无释放，重现了女士报告的现象。
\end{enumerate}

% ============================================================================
% 第二部分  阅读
% ============================================================================
\GaokaoPart{第二部分}{阅读（每小题 2.5 分，满分 50 分）}

\begin{center}
\renewcommand{\arraystretch}{1.35}
\begin{tabular}{|c|c|c|c|c|c|c|c|c|c|c|}
\hline
\textbf{题号} & 11 & 12 & 13 & 14 & 15 & 16 & 17 & 18 & 19 & 20 \\
\hline
\textbf{答案} & \textbf{C} & \textbf{A} & \textbf{D} & \textbf{B} & \textbf{C} & \textbf{A} & \textbf{D} & \textbf{B} & \textbf{A} & \textbf{C} \\
\hline
\textbf{题号} & 21 & 22 & 23 & 24 & 25 & 26 & 27 & 28 & 29 & 30 \\
\hline
\textbf{答案} & \textbf{B} & \textbf{D} & \textbf{C} & \textbf{A} & \textbf{B} & \textbf{F} & \textbf{C} & \textbf{A} & \textbf{G} & \textbf{D} \\
\hline
\end{tabular}
\end{center}

\begin{enumerate}[label=\textbf{\arabic*.}, start=11, leftmargin=2.2em, itemsep=2pt, topsep=3pt]
  \item \textbf{C}\quad 第三行建议保留并行设置，记录各进程使用的资源名排查重名。
  \item \textbf{A}\quad 保留失败时的次序并移除其余用例，可检验这一对能否单独重现故障。
  \item \textbf{D}\quad 用户只报告了单跑和整组的结果，尚未确定“那一对”中的前序用例。
  \item \textbf{B}\quad 开幕当天展厅无网，馆员要求按现场条件试用。
  \item \textbf{C}\quad 操场原图刚在联网时打开过，展机留有副本。
  \item \textbf{A}\quad 清除留存副本后，断网查看原图可检验展览包自身的内容。
  \item \textbf{D}\quad 食堂原图在第一次断网时显示为灰框，馆员先复查这一项。
  \item \textbf{B}\quad 处理函数在 250\,ms 进入、262\,ms 返回，相隔 12 毫秒。
  \item \textbf{A}\quad 从函数返回至客户端读完响应体，用时为 $310 - 262 = 48$ 毫秒。
  \item \textbf{C}\quad 新字段名 \texttt{handler\_duration\_ms} 指明计时对象为处理函数。
  \item \textbf{B}\quad 开始申请与拿到连接的时刻均缺，只补一端仍无法相减求得等待时长。
  \item \textbf{D}\quad 早期脚本将设备实测为零与未获应答均记为零。
  \item \textbf{C}\quad 所存时刻仅在读取成功时更新，记录最近一次成功取得读数的时刻。
  \item \textbf{A}\quad 所查提交记录能确定 A 分支出现该段文字早两天，代码迁移方向仍未确定。
  \item \textbf{B}\quad 前者对应设备报告零字节，后者对应本轮未取得读数。
  \item \textbf{F}\quad 两段校验内容与次序相同，“这两处调用”承接两个入口对共同函数的调用。
  \item \textbf{C}\quad 旧客户端依赖原错误码；B 项统一两端错误码会改变既有行为。
  \item \textbf{A}\quad 锁定次数已有单独议题核对，本次改动说明仅附链接。
  \item \textbf{G}\quad 两个入口各留一份完整响应，与下句“这两份响应”直接衔接；E 项新增与需求无关的落盘摘要。
  \item \textbf{D}\quad 摘要交代已提取重复校验并核对响应，末尾附上锁定规则议题链接。
\end{enumerate}

% ============================================================================
% 第三部分  语言运用
% ============================================================================
\GaokaoPart{第三部分}{语言运用（共两节，满分 30 分）}

\GaokaoSection{第一节：完形填空}{（每小题 1 分，满分 15 分）}

\begin{center}
\setlength{\tabcolsep}{4.8pt}
\renewcommand{\arraystretch}{1.35}
\begin{tabular}{|c|c|c|c|c|c|c|c|c|c|c|c|c|c|c|c|}
\hline
\textbf{题号} & 31 & 32 & 33 & 34 & 35 & 36 & 37 & 38 & 39 & 40 & 41 & 42 & 43 & 44 & 45 \\
\hline
\textbf{答案} & \textbf{B} & \textbf{D} & \textbf{A} & \textbf{C} & \textbf{B} & \textbf{A} & \textbf{D} & \textbf{C} & \textbf{B} & \textbf{A} & \textbf{D} & \textbf{C} & \textbf{A} & \textbf{D} & \textbf{B} \\
\hline
\end{tabular}
\end{center}

\begin{enumerate}[label=\textbf{\arabic*.}, start=31, leftmargin=2.2em, itemsep=2pt, topsep=3pt]
  \item \textbf{B}\quad 提交号哈希值相同，说明两份运行记录指向同一次提交。
  \item \textbf{D}\quad 函数收到的是带时区偏移量的完整时间字符串。
  \item \textbf{A}\quad 两份断言写的都是“导出”，即期望结果相同。
  \item \textbf{C}\quad \texttt{UTC} 周日 16:15 即北京时间周一 00:15，两行记录的是同一个时刻。
  \item \textbf{B}\quad 两行日志分别写着 11 月 30 日与 11 月 29 日，相差一天的是当地日期。
  \item \textbf{A}\quad 旧函数先把输入换算成机器本地时间，再从中取星期。
  \item \textbf{D}\quad 开发机保持代码、输入与断言不变，只把测试进程的时区调成 \texttt{UTC}。
  \item \textbf{C}\quad 本地出现了与 CI 相同的报错，失败已复现。
  \item \textbf{B}\quad 改动进程时区后，同一条测试的导出判断随之改变。
  \item \textbf{A}\quad 若只调 CI 而不改函数，换台机器仍要看运行环境的时区设置。
  \item \textbf{D}\quad 显式接收业务时区并据此判断周几的，是导出判断函数。
  \item \textbf{C}\quad 同一条用例分别在 \texttt{UTC+8} 与 \texttt{UTC} 两种进程时区下运行。
  \item \textbf{A}\quad 北京时间周一 00:00 前一分钟为周日 23:59。
  \item \textbf{D}\quad 周日 23:59 按北京时间仍属周日，按规则应当不导出。
  \item \textbf{B}\quad 实测覆盖的是周界两侧各一条用例，两条在两种设置下均符合规则。
\end{enumerate}

\GaokaoSection{第二节：语法填空}{（每小题 1.5 分，满分 15 分）}

\noindent{\small 每空答对得 1.5 分，答错或未答得 0 分。无提示词空按句法及原记录的时间、条件和确认范围判分，扩大确认范围的记 0 分。}

\begin{center}
\renewcommand{\arraystretch}{1.28}
\begin{tabularx}{\linewidth}{|c|c|X|}
\hline
\textbf{题号} & \textbf{标准答案} & \textbf{评分说明} \\
\hline
46 & \textbf{把} & 引出回退的两台实例。填“将”同等给分。 \\
\hline
47 & \textbf{从} & 新版本是回退起点。填“由”同等给分。 \\
\hline
48 & \textbf{到} & 旧版本是回退终点。填“至”同等给分。 \\
\hline
49 & \textbf{已暂停} & 14:05 已停止向其余实例继续发布。填“待暂停”“未暂停”或原词直填记 0 分。 \\
\hline
50 & \textbf{只} & 本轮回滚限于已升级的两台实例。填“仅”同等给分；扩大范围记 0 分。 \\
\hline
51 & \textbf{已合入} & 修复提交已进入主分支，合并单由“待”转“已”。填“待合入”或原词直填记 0 分。 \\
\hline
52 & \textbf{过} & “有没有……重新测过”询问交班前是否有过该次测试。填“了”记 0 分。 \\
\hline
53 & \textbf{复测} & 与“初测”相对，表示再次测试。填“重测”同等给分。 \\
\hline
54 & \textbf{等} & 先完成部署，再用线上订单核对。填“待”同等给分。 \\
\hline
55 & \textbf{才} & 与“只有”呼应，标明填写“生产验证”一栏的必要条件。填“方”同等给分。 \\
\hline
\end{tabularx}
\end{center}

% ============================================================================
% 第四部分  写作
% ============================================================================
\GaokaoPart{第四部分}{写作（共两节，满分 40 分）}

\GaokaoSection{第一节：应用文写作}{（第 56 题，满分 15 分）}

\noindent\textbf{【参考作答一】}

\begin{quote}
\vspace{-2pt}
\kaishu\noindent
“张工，收到。这周末我排不开。旧版移动端还靠 \texttt{total} 显示订单金额，直接删会影响显示。周一上午我先整理一版兼容方案发您，上线时间等方案定了再约。”\\[1pt]
如果周一只能先碰方案，把“上午我先整理一版兼容方案发您”换成“我们先碰一下兼容方案”，先把兼容边界钉住，别顺手把交付时间也钉死。
\vspace{-2pt}
\end{quote}

\noindent{\small\textbf{评注：}“上午我先整理一版兼容方案发您”由助理添入；按末句说明替换后，发出的是“周一我们先碰一下兼容方案”，李华只答应了这一步。考生直接采用替换后的表述（见参考作答二），同等给分。}

\smallskip
\needspace{5\baselineskip}
\noindent\textbf{【参考作答二】}

\begin{quote}
\vspace{-2pt}
\kaishu\noindent
“张工，收到。这周末我排不开。旧版移动端还在读 \texttt{total} 展示订单金额，现在删掉，旧版就取不到这个金额了。周一我可以跟您对一下兼容方案，看看旧版怎么过渡；上线时间等方案商量好再定。”\\[1pt]
引号里这段可以直接发给张工；周一具体几点聊，等你们约好再补。
\vspace{-2pt}
\end{quote}

\noindent\textbf{评分项目与五档标准：}
\begin{enumerate}[label=（\arabic*）, leftmargin=2.8em, labelsep=0pt, itemsep=0.5pt, topsep=1.5pt]
  \item \textbf{字段依赖（5 分）：}指明旧版移动端（1 分），写明其仍读取 \texttt{total}（2 分），说明删除会影响订单金额展示（2 分）。
  \item \textbf{时间安排（4 分）：}周末排不开（1 分），周一可商谈方案（1 分），上线时间待定（2 分）。附替换说明的按替换后的文本评定。
  \item \textbf{发送与说明（4 分）：}正文可单独复制发送（2 分），说明指明发送范围或具体替换位置（2 分）。
  \item \textbf{对象与口吻（2 分）：}给张工的消息与给李华的说明区分清楚（1 分），口吻得当（1 分）。
\end{enumerate}

\noindent{\small 按总分定档：第五档（13--15 分）事实准确，安排清楚，可按说明直接发送；第四档（10--12 分）主要事实齐全，个别时间措辞或说明略欠明确；第三档（7--9 分）能回应请求，遗漏字段影响或使用说明；第二档（4--6 分）遗漏关键依赖、混淆对象或擅定交付时间；第一档（1--3 分）仅有零散相关表述。}

\GaokaoSection{第二节：读后续写}{（第 57 题，满分 25 分）}

\noindent\textbf{【官方参考范文】}

\begin{quote}
\vspace{-2pt}
\noindent\textbf{Paragraph 1:}\par
\kaishu\noindent\hspace*{2em}我请她先撤掉那三个空格，把原稿按日志里的位置分成两块。分两次送进去，“桥”坏了；把同样的字节合起来再送一次，字就回来了。我翻到读取循环，每读一块，程序都当作一段完整文字解码。读到第一块末尾，它已经把那一个字节换成了替换符，后两个字节来时，前一个早已交给了排版程序。\par
\smallskip
\noindent\textbf{Paragraph 2:}\par
\kaishu\noindent\hspace*{2em}校样送印前，编辑又发来那份没有加空格的稿子。我改用能跨块接着读的解码器，把“桥”分别切在第一个和第二个字节之后，逐字核对，两份预览都与原稿相同。她看过校样，点了送印。那份原稿后来进了测试目录，旁边只附了一行说明：分两次读。
\vspace{-2pt}
\end{quote}

\noindent\textbf{评分项目与五档标准：}
\begin{enumerate}[label=（\arabic*）, leftmargin=2.8em, labelsep=0pt, itemsep=0.5pt, topsep=1.5pt]
  \item \textbf{情节承接（5 分）：}接续分块复现试验与送印前复核，两段先后顺序及人物所知与段首连贯。
  \item \textbf{故障判断（7 分）：}写出可区分原因的排查动作（3 分），结合字节完整、分块后出错及空格平移坏字位置的现象说明根因（4 分）。
  \item \textbf{处理与复验（8 分）：}处理方法能保全跨块多字节字符（4 分），并用原稿或跨块切分边界样例核对预览结果（4 分，其他合理修复方案同等给分）。
  \item \textbf{叙述与表达（5 分）：}技术细节与因果关系准确（3 分），语体克制自然，结尾有前文细节可据（2 分）。
\end{enumerate}

\begin{itemize}[leftmargin=1.8em, itemsep=0.5pt, topsep=1.5pt]
  \item \textbf{第五档（21--25 分）：}两段完整连贯；排查动作有辨别力，根因判断、跨块处理与边界复验相互印证；叙述平实克制，收束有前文可据。
  \item \textbf{第四档（16--20 分）：}根因判断合理，处理与复核经过清楚；个别复验细节略简，但整体因果闭环完整。
  \item \textbf{第三档（11--15 分）：}能写出分块读取与乱码的关系，处理方法或复验过程交代不足。
  \item \textbf{第二档（6--10 分）：}叙事与材料有关，仍停留在重试、手动增删空格等权宜处理，根因交代不清。
  \item \textbf{第一档（1--5 分）：}仅有零散叙述，未能有效续接所给段首句。
\end{itemize}

\end{document}
